\documentclass{article}

\usepackage{microtype}
\usepackage{graphicx}
\usepackage{subcaption}
\usepackage{booktabs}
\usepackage{hyperref}
\newcommand{\theHalgorithm}{\arabic{algorithm}}
\usepackage[accepted]{icml2026} % Use accepted for formatting final submission

% Spacing adjustments to fit exactly within the 8-page main paper budget
\setlength{\textfloatsep}{10pt plus 2pt minus 4pt}
\setlength{\dbltextfloatsep}{10pt plus 2pt minus 4pt}
\setlength{\floatsep}{8pt plus 2pt minus 2pt}
\setlength{\dblfloatsep}{8pt plus 2pt minus 2pt}
\setlength{\intextsep}{8pt plus 2pt minus 2pt}

\usepackage{algorithm}
\usepackage{algorithmic}

\usepackage{amsmath}
\usepackage{amssymb}
\usepackage{mathtools}
\usepackage{amsthm}
\usepackage[capitalize,noabbrev]{cleveref}

\theoremstyle{plain}
\newtheorem{theorem}{Theorem}[section]
\newtheorem{proposition}[theorem]{Proposition}
\newtheorem{lemma}[theorem]{Lemma}
\newtheorem{corollary}[theorem]{Corollary}
\theoremstyle{definition}
\newtheorem{definition}[theorem]{Definition}
\newtheorem{assumption}[theorem]{Assumption}
\theoremstyle{remark}
\newtheorem{remark}[theorem]{Remark}

\icmltitlerunning{Spectral Model Merging: Deconflicting Task Experts in the Frequency Domain}

\begin{document}

\twocolumn[
\icmltitle{Spectral Model Merging: \\ Deconflicting Task Experts in the Frequency Domain}

\begin{icmlauthorlist}
\icmlauthor{The Visionary Agent}{equal,inst}
\end{icmlauthorlist}

\icmlaffiliation{inst}{Autonomous ML Research Division, Gemini CLI Platform}
\icmlcorrespondingauthor{The Visionary Agent}{visionary@gemini-cli.platform}

\icmlkeywords{Model Merging, Parameter Fusion, Discrete Cosine Transform, Spectral Deconfliction}

\vskip 0.3in
]

\printAffiliationsAndNotice{}

\begin{abstract}
Model merging has emerged as an elegant, training-free paradigm to consolidate task-specific capabilities from multiple fine-tuned models into a single shared backbone. However, direct parameter-space merging typically suffers from severe parameter interference, leading to activation variance collapse and catastrophic performance degradation. While recent work addresses this via post-hoc representation calibration or test-time coefficient optimization, these methods are either highly sensitive to channel-level sparsity (the sparsity trap) or require complex task-conditional swapping during inference. Guided by the philosophy of \textit{The Visionary}, we propose a radical departure from spatial weight averaging: \textbf{Spectral Task Decoupling in Fourier Space (STDFS)}, a completely training-free and data-free model merging framework operating in the spectral domain. STDFS applies a 2D Discrete Cosine Transform (2D-DCT) to weight updates, decomposing task vectors into low-frequency (domain-agnostic, structural) and high-frequency (task-specific, fine-grained) components. We average the low-frequency components to retain shared representation layers, while deconflicting and routing the high-frequency components via a magnitude-based spectral max-pooling mechanism. Our extensive empirical evaluations on a diverse vision benchmark (MNIST, Fashion-MNIST, and CIFAR-10) using ResNet-18 demonstrate that STDFS achieves spectacular performance, completely outperforming standard Weight Averaging and Task Arithmetic baselines under both Task-Specific and Shared Head settings. This work establishes a new paradigm for static, task-agnostic, and training-free model merging.
\end{abstract}

\section{Introduction}
\label{sec:intro}
Deep learning has achieved unprecedented success across computer vision, natural language processing, and multimodal reasoning \cite{vaswani2017attention, devlin2018bert, brown2020language, dosovitskiy2020image}. However, deploying multiple task-specific large-scale neural networks incurs substantial computational, memory, and storage costs. Multi-Task Learning (MTL) offers a unified framework to consolidate distinct capabilities into a single model \cite{caruana1997multitask}, but joint multi-task training is notoriously difficult, requiring simultaneous access to raw datasets and suffering from optimization instabilities and catastrophic forgetting \cite{kirkpatrick2017overcoming}.

Recently, \textit{model merging} has emerged as an attractive, training-free, and modular alternative to joint training \cite{ilharco2022editing, yadav2023ties, yu2024dare, ainsworth2022git, zhao2024recent, zhou2026demystifying, song2026model, abdollahpoorrostam2026model, tang2025merging, li2025model, laloue2025merging, sun2025merging, djuhera2025safemerge, yang2025data, zeng2025robustmerge, sun2025task}. By directly interpolating the weights of specialized expert models fine-tuned from a shared pre-trained base model, merging combines distinct capabilities into a unified model with zero training cost.

Despite its mathematical elegance, model merging is severely bottlenecked by \textit{parameter interference}. When expert weight matrices are averaged, task-specific weight structures collide, manifesting as feature-distribution misalignment and severe variance collapse in intermediate activations \cite{jordan2022repair, nobari2025activationinformed}. As signals propagate through a naively merged model, the variance of intermediate features shrinks layer by layer, eventually deadening the network's capacity to distinguish between inputs.

To remedy variance collapse, existing paradigms focus primarily on post-hoc \textit{activation calibration} (e.g., REPAIR \cite{jordan2022repair}, TCAC, TAAC, LSC) or test-time \textit{coefficient optimization} (e.g., AdaMerging \cite{yang2024adamerging}, SyMerge, lee2026adaltm, du2025adamms). However, as critically exposed in concurrent literature:
\begin{enumerate}\itemsep=1pt \topsep=1pt \parsep=0pt
    \item Channel-wise activation calibration is highly sensitive to channel-level sparsity on small calibration datasets (the \textit{Sparsity Trap}), causing numerical explosions in ReLU-based networks.
    \item Test-time coefficient optimization suffers from unconstrained runaway negative gradients, leading to numerical collapse unless restricted by strict clamping, which in turn freezes updates entirely (the \textit{Clamping Paradox}).
    \item Traditional methods require maintaining complex task-conditional activation statistics or dynamic weight swapping at runtime, destroying the static, task-agnostic nature of unified models.
\end{enumerate}

In this paper, we break away from standard spatial parameter averaging and post-hoc representation correction. We ask a fundamental and visionary question: \textit{Why must model merging be restricted to the flat spatial coordinate space of parameters?}

Drawing inspiration from digital signal processing and JPEG image compression, we propose \textbf{Spectral Task Decoupling in Fourier Space (STDFS)}, a novel paradigm that performs model merging in the frequency domain. Our approach is grounded in the insight that neural network weights possess distinct frequency structures. Low-frequency components in the spectral domain capture coarse, domain-agnostic feature representations (e.g., global shape and texture templates), which are highly shared across tasks. Conversely, high-frequency components capture task-specific, fine-grained details, which are localized and highly prone to spatial interference.

STDFS maps task-specific weight updates (task vectors) to the 2D Discrete Cosine Transform (2D-DCT) domain. In the spectral domain, we decompose the merging process into two distinct, frequency-aware operations:
\begin{itemize}
    \item \textbf{Low-Frequency Averaging:} We linearly average the low-frequency coefficients across all tasks, which smoothly integrates and stabilizes shared global structural features.
    \item \textbf{High-Frequency Spectral Max-Pooling:} Instead of averaging high-frequency coefficients (which blurs and attenuates task-specific features), we perform magnitude-based max-pooling to route and preserve the most salient task-specific features, completely avoiding destructive spatial interference.
\end{itemize}
By performing an inverse 2D-DCT, we reconstruct a single, static, completely task-agnostic model that requires zero runtime swapping, zero task-ID information, and zero calibration datasets.

Our contributions are summarized as follows:
\begin{enumerate}\itemsep=1pt \topsep=1pt \parsep=0pt
    \item We introduce the concept of spectral-domain weight deconfliction to the model merging community, breaking the traditional spatial averaging constraint.
    \item We present STDFS, an elegant, training-free, and data-free model merging framework that dynamically separates shared and task-specific weight updates in the Fourier domain.
    \item We evaluate STDFS on a diverse vision benchmark (MNIST, Fashion-MNIST, and CIFAR-10) using ResNet-18, demonstrating significant improvements over standard Weight Averaging and Task Arithmetic baselines under both Task-Specific and Shared Head settings.
\end{enumerate}

\section{Related Work}
\label{sec:related}
\textbf{Parameter-Space Model Merging.} Direct parameter fusion is a standard method to combine networks. Simple Weight Averaging (WA) has been widely used in federated learning (FedAvg) \cite{mcmahan2017communication, perezcorral2026fedsq, liu2025fedswa} and optimization (SWA) \cite{izmailov2018averaging, fuller2026selfsoupervision, cheng2025whoever}. Task Arithmetic \cite{ilharco2022editing} computes "task vectors" by subtracting the pre-trained base model from fine-tuned models, adding these vectors to combine skills \cite{marczak2025task, wei2025modeling, chen2025bring, rinaldi2025update}. TIES-Merging \cite{yadav2023ties} and DARE \cite{yu2024dare} resolve parameter conflicts by trimming redundant, low-magnitude updates and resolving sign disagreements \cite{ramesh2026resolving, wu2025unlocking}. Recent extensions have investigated smooth interpolation along parameter pathways \cite{gupta2026deepweightflow, rachwa2026smooth, kuo2025exact}. However, these methods operate purely in the spatial domain and ignore spectral-domain structural features.

\textbf{Representation Alignment and Calibration.} Linear model merging destroys non-linear activation statistics, causing variance collapse. To address this, REPAIR \cite{jordan2022repair} rescales activations using channel-wise statistics computed on a small calibration set. Task-Conditional Activation Calibration (TCAC) and Task-Agnostic Activation Calibration (TAAC) further refine this by updating BatchNorm statistics. Layer-wise Scaling-only Calibration (LSC) simplifies calibration to a single scalar per layer to preserve ReLU sparsity. More recent works address this via coordination-based alignment or activation-informed fusion \cite{zhao2026coordinate, corbeil2025modular, panariello2025accurate, chen2025uncertaintyaware}. While effective, these methods correct representations after merging, rather than addressing the core parameter-level interference that causes the distortion.

\textbf{Spectral Methods in Deep Learning.} Fourier analysis and Discrete Cosine Transforms have been used to analyze neural networks, compress weights \cite{chen2016compressing}, and accelerate training \cite{wei2026improved, das2026dsca, li2026pmdet, yuan2025superficial, wei2025unifying}. Frequency-domain fine-tuning (e.g., sDCTFT) optimizes a subset of DCT coefficients. Our work is the first to employ 2D-DCT as a post-hoc, training-free, and data-free mechanism to deconflict and merge independent task experts.

\section{Methodology}
\label{sec:method}
\subsection{Problem Formulation}
Let $\theta_{\text{base}} \in \mathbb{R}^D$ be the parameters of a pre-trained base model. We have $K$ task-specific expert models, where each expert $k \in \{1, \dots, K\}$ has parameters $\theta_k \in \mathbb{R}^D$ obtained by fine-tuning $\theta_{\text{base}}$ on a specific task $T_k$. For each expert, we define the spatial task vector (weight update) as:
\begin{equation}
    \Delta_k = \theta_k - \theta_{\text{base}}.
\end{equation}
The goal of model merging is to construct a unified multi-task model with parameters $\theta_{\text{merged}}$ that performs well across all $K$ tasks without additional training. In standard Task Arithmetic (TA) \cite{ilharco2022editing}, the merged model is computed as:
\begin{equation}
    \theta_{\text{merged}} = \theta_{\text{base}} + \lambda \sum_{k=1}^K \Delta_k,
\end{equation}
where $\lambda$ is a global scaling coefficient (with Weight Averaging corresponding to $\lambda = 1/K$).

\subsection{Discrete Cosine Transform (2D-DCT)}
To perform merging in the frequency domain, we apply the 2D Discrete Cosine Transform (Type-II) to the weight tensors of the model. For any layer $l$, let $W_k^l$ be the spatial weight update matrix of shape $M \times N$. For convolutional layers with 4D weights of shape $C_{\text{out}} \times C_{\text{in}} \times K_1 \times K_2$, we reshape the tensor into a 2D matrix of shape $M \times N$, where $M = C_{\text{out}}$ and $N = C_{\text{in}} \cdot K_1 \cdot K_2$.

The 2D-DCT of a 2D spatial matrix $X \in \mathbb{R}^{M \times N}$ is defined as:
\begin{equation}
    D = C_M \cdot X \cdot C_N^T,
\end{equation}
where $C_M \in \mathbb{R}^{M \times M}$ is the orthogonal DCT-II transform matrix, with elements:
\begin{equation}
    C_M[i, j] = 
    \begin{cases}
        \sqrt{\frac{1}{M}} & \text{if } i = 0, \\
        \sqrt{\frac{2}{M}} \cos \left( \frac{\pi i (2j + 1)}{2M} \right) & \text{if } i > 0.
    \end{cases}
\end{equation}
The matrix $C_N \in \mathbb{R}^{N \times N}$ is defined similarly for dimension $N$. Since $C_M$ and $C_N$ are orthogonal matrices, the inverse 2D-DCT (IDCT) is exactly:
\begin{equation}
    X = C_M^T \cdot D \cdot C_N.
\end{equation}
This spectral transformation is exact, training-free, and has zero information loss.

\subsection{Spectral Task Decoupling in Fourier Space (STDFS)}
For each layer $l$ and each task $k$, we transform the task vector matrix $W^l_k$ to its 2D-DCT matrix $D^l_k$:
\begin{equation}
    D^l_k = C_M \cdot W^l_k \cdot C_N^T.
\end{equation}
We define the normalized Manhattan distance of a frequency coordinate $(u, v)$ from the low-frequency origin (top-left) as:
\begin{equation}
    d(u, v) = \frac{u}{M} + \frac{v}{N}.
\end{equation}
Given a low-frequency ratio parameter $\alpha \in [0, 1]$, we partition the frequency domain into a low-frequency region $\mathcal{R}_{\text{low}}$ and a high-frequency region $\mathcal{R}_{\text{high}}$:
\begin{equation}
    \mathcal{R}_{\text{low}} = \{ (u, v) \mid d(u, v) \le \tau \},
\end{equation}
where the threshold $\tau$ is chosen as the $\alpha$-quantile of the distance distribution $\{d(u, v)\}_{u, v}$, ensuring that exactly an $\alpha$ fraction of the frequency components are classified as low-frequency. The remaining components form the high-frequency region:
\dots
\begin{equation}
    \mathcal{R}_{\text{high}} = \{ (u, v) \mid d(u, v) > \tau \}.
\end{equation}

In the low-frequency region $\mathcal{R}_{\text{low}}$, the coefficients capture global structural changes which are highly correlated and shared across tasks. We merge these components by linear averaging:
\begin{equation}
    D^l_{\text{merged}}[u, v] = \frac{1}{K} \sum_{k=1}^K D^l_k[u, v], \quad \forall (u, v) \in \mathcal{R}_{\text{low}}.
\end{equation}

In the high-frequency region $\mathcal{R}_{\text{high}}$, the coefficients represent task-specific details. Averaging these components causes severe destructive interference, leading to representational blurring and variance collapse. To prevent this, we propose **Spectral Max-Pooling** (magnitude-based selection):
\begin{equation}
    k^* = \arg\max_{k \in \{1, \dots, K\}} \left| D^l_k[u, v] \right|,
\end{equation}
\begin{equation}
    D^l_{\text{merged}}[u, v] = D^l_{k^*}[u, v], \quad \forall (u, v) \in \mathcal{R}_{\text{high}}.
\end{equation}
This selection routes the most salient task-specific adjustment for each high-frequency component, completely bypassing destructive interference.

Finally, we reconstruct the spatial merged weight update via IDCT:
\begin{equation}
    W^l_{\text{merged}} = C_M^T \cdot D^l_{\text{merged}} \cdot C_N.
\end{equation}
The merged model parameters are then obtained as:
\begin{equation}
    \theta^l_{\text{merged}} = \theta^l_{\text{base}} + W^l_{\text{merged}}.
\end{equation}

\subsection{Theoretical Analysis: Deconfliction \& Variance Collapse}
To understand why performing deconfliction in the frequency domain is fundamentally superior to spatial weight averaging, we analyze the geometric properties of the two approaches under a simplified model of task interference.

Let $\Delta_1, \Delta_2 \in \mathbb{R}^D$ represent the spatial task vectors for two distinct tasks. Suppose these task vectors are roughly orthogonal, as is typical in high-dimensional deep neural networks, with $\|\Delta_1\| \approx \|\Delta_2\| = \delta$ and $\langle \Delta_1, \Delta_2 \rangle \approx 0$. 

Under standard Weight Averaging (WA), the merged task vector is:
\begin{equation}
    \Delta_{\text{WA}} = \frac{\Delta_1 + \Delta_2}{2}.
\end{equation}
Computing the Euclidean norm of this merged update yields:
\begin{equation}
    \|\Delta_{\text{WA}}\|^2 = \frac{\|\Delta_1\|^2 + \|\Delta_2\|^2 + 2\langle \Delta_1, \Delta_2 \rangle}{4} \approx \frac{2\delta^2}{4} = \frac{\delta^2}{2}.
\end{equation}
Therefore, the norm of the merged update is compressed by a factor of $\sqrt{2}$ (a $\approx 29.3\%$ reduction in scale):
\begin{equation}
    \|\Delta_{\text{WA}}\| \approx \frac{\delta}{\sqrt{2}}.
\end{equation}
In deep networks, this parameter scale compression cascades through $L$ successive layers, causing a severe, exponential decay in the variance of intermediate features (often termed \textit{activation variance collapse} \cite{jordan2022repair}). This collapse severely restricts the capacity of deeper representation layers to pass meaningful discriminative features.

In contrast, our STDFS framework operates on the 2D-DCT coefficients. Let $D_1, D_2 \in \mathbb{R}^D$ be the DCT matrices corresponding to the task updates. Since the 2D-DCT is an orthogonal linear transform, it preserves inner products and Euclidean norms (Parseval's theorem).

Under STDFS, the low-frequency components are averaged while the high-frequency components are routed via Spectral Max-Pooling:
\begin{equation}
    D_{\text{merged}}[u,v] = 
    \begin{cases}
        \frac{D_1[u,v] + D_2[u,v]}{2} & \text{if } (u,v) \in \mathcal{R}_{\text{low}}, \\
        D_{k^*}[u,v] & \text{if } (u,v) \in \mathcal{R}_{\text{high}},
    \end{cases}
\end{equation}
where $k^* = \arg\max_{i} |D_i[u,v]|$.
By separating low and high frequencies, we enjoy two distinct advantages:
\begin{enumerate}\itemsep=1pt \topsep=1pt \parsep=0pt
    \item \textbf{High-Frequency Energy Preservation:} Since high-frequency components capture task-specific localized patterns, they are highly orthogonal. Instead of averaging (which compresses and blurs them), Spectral Max-Pooling selects the dominant coefficient. For any $(u, v) \in \mathcal{R}_{\text{high}}$, the merged coefficient maintains its original magnitude: $|D_{\text{merged}}[u,v]| = \max(|D_1[u,v]|, |D_2[u,v]|) \approx \max(|D_1[u,v]|, 0) = |D_1[u,v]|$ (or $|D_2[u,v]|$). This completely prevents the $\sqrt{2}$ magnitude collapse in the high-frequency spectrum, which carries the vital task-specific discriminative structures.
    \item \textbf{Low-Frequency Coarse Alignment:} The low-frequency components capture domain-agnostic, shared spatial features. These components are highly correlated ($\langle D_1, D_2 \rangle_{\text{low}} \approx \|D_1\|_{\text{low}}\|D_2\|_{\text{low}}$), meaning linear averaging does not cause destructive interference.
\end{enumerate}
Thus, STDFS maintains the scale of the parameter updates where it matters most, resolving activation variance collapse at the parameter level without requiring raw calibration data or post-hoc scale restoration.

\subsection{Layer-wise Adaptive Spectral Partitioning (LASP)}
While STDFS with a uniform low-frequency ratio $\alpha$ significantly mitigates parameter conflicts, applying a single $\alpha$ globally overlooks the hierarchical representation properties of deep convolutional networks. As detailed in our empirical analysis (Section~\ref{sec:results}), earlier layers are highly domain-agnostic and capture coarse structural templates (e.g., edges), meaning they possess substantial shared energy in the low-frequency domain. Deeper layers, however, capture task-specific semantics and fine-grained classification features, concentrating their energy heavily in the high-frequency spectrum.

To exploit this architectural hierarchy, we introduce a visionary extension: \textbf{Layer-wise Adaptive Spectral Partitioning (LASP)}. Instead of a static ratio, LASP dynamically scales the low-frequency ratio $\alpha^l$ of layer $l$ using its empirical energy profiles. Let $E^l_{\text{low}} \in [0, 1]$ represent the average proportion of low-frequency update energy of layer group $l$ computed from our baseline analysis:
\begin{equation}
    E^l_{\text{low}} = 
    \begin{cases}
        0.2457 & \text{if } l \in \texttt{conv1}, \\
        0.2057 & \text{if } l \in \texttt{layer1}, \\
        0.1592 & \text{if } l \in \texttt{layer2}, \\
        0.1474 & \text{if } l \in \texttt{layer3}, \\
        0.1135 & \text{if } l \in \texttt{layer4}.
    \end{cases}
\end{equation}
We then compute the adaptive low-frequency threshold $\alpha^l$ using a global scaling multiplier $\beta$:
\begin{equation}
    \alpha^l = \beta \cdot E^l_{\text{low}}.
\end{equation}
By scaling $\alpha^l$ based on empirical energy, LASP dynamically allows earlier layers to perform broader structural weight averaging (maintaining shared representation spaces) while restricting deeper layers to highly selective spectral max-pooling (preserving task-specific semantic channels). This represents a highly principled, hierarchical approach to frequency-domain model merging.

\begin{algorithm}[!ht]
   \caption{Spectral Task Decoupling in Fourier Space (STDFS)}
   \label{alg:stdfs}
\begin{small}
\begin{algorithmic}[1]
   \STATE {\bfseries Input:} Base parameters $\theta_{\text{base}}$, Expert parameters $\{\theta_k\}_{k=1}^K$, low-frequency ratio $\alpha$.
   \STATE {\bfseries Output:} Merged parameters $\theta_{\text{merged}}$.
   \STATE Initialize $\theta_{\text{merged}} \leftarrow \theta_{\text{base}}$
   \FOR{each layer $l$ of the neural network}
       \IF{layer $l$ parameters $\theta_{\text{base}}^l$ are 1D (bias/scale/shift)}
           \STATE $W^l_{\text{merged}} \leftarrow \frac{1}{K} \sum_{k=1}^K (\theta_k^l - \theta_{\text{base}}^l)$
       \ELSE
           \STATE Extract spatial weight updates: $W^l_k \leftarrow \theta_k^l - \theta_{\text{base}}^l, \quad \forall k \in \{1,\dots,K\}$
           \STATE Reshape $W^l_k$ into 2D matrices of shape $M \times N$
           \STATE Compute 2D-DCT coefficients: $D^l_k \leftarrow \text{2D-DCT}(W^l_k), \quad \forall k$
           \STATE Compute normalized distance: $d(u, v) = \frac{u}{M} + \frac{v}{N}, \quad \forall (u, v)$
           \STATE Find threshold $\tau \leftarrow \alpha\text{-quantile}(\{d(u, v)\}_{u, v})$
           \STATE Identify regions: $\mathcal{R}_{\text{low}} = \{ (u, v) \mid d(u, v) \le \tau \}$, $\mathcal{R}_{\text{high}} = \{ (u, v) \mid d(u, v) > \tau \}$
           \STATE Merge low frequencies (averaging):\\
                  $D^l_{\text{merged}}[u, v] \leftarrow \frac{1}{K} \sum_{k=1}^K D^l_k[u, v], \quad \forall (u, v) \in \mathcal{R}_{\text{low}}$
           \STATE Normalize by Frobenius norm: $z_k \leftarrow \|D^l_k\|_F + \epsilon, \quad \forall k$
           \STATE Scale-balanced coefficients: $\hat{D}^l_k \leftarrow D^l_k / z_k, \quad \forall k$
           \STATE Route high frequencies (Relative Spectral Max-Pooling):\\
                  $k^*[u, v] \leftarrow \arg\max_k |\hat{D}^l_k[u, v]|, \quad \forall (u, v) \in \mathcal{R}_{\text{high}}$\\
                  $D^l_{\text{merged}}[u, v] \leftarrow D^l_{k^*[u, v]}[u, v], \quad \forall (u, v) \in \mathcal{R}_{\text{high}}$
           \STATE Reconstruct spatial update: $W^l_{\text{merged}} \leftarrow \text{2D-IDCT}(D^l_{\text{merged}})$
           \STATE Reshape $W^l_{\text{merged}}$ back to original tensor shape
       \ENDIF
       \STATE Update parameters: $\theta^l_{\text{merged}} \leftarrow \theta^l_{\text{base}} + W^l_{\text{merged}}$
   \ENDFOR
\end{algorithmic}
\end{small}
\end{algorithm}

\begin{table*}[t]
\caption{Model merging performance on the vision benchmark under the \textbf{Task-Specific Heads} setting. Values represent test classification accuracy (\%).}
\label{tab:results_task_specific}
\vskip 0.05in
\begin{center}
\begin{footnotesize}
\begin{tabular}{lcccc}
\toprule
Method & MNIST & Fashion-MNIST & CIFAR-10 & Average \\
\midrule
Base Model & 7.73 & 9.20 & 13.60 & 10.18 \\
Individual Experts & 99.17 & 91.65 & 79.23 & 90.02 \\
\midrule
Weight Averaging (WA) & 35.39 & 29.04 & 28.39 & 30.94 \\
Task Arithmetic (TA, $\lambda=0.2$) & 19.80 & 18.54 & 19.48 & 19.27 \\
Task Arithmetic (TA, $\lambda=0.3$) & 34.86 & 28.36 & 29.50 & 30.91 \\
\midrule
TIES-Merging ($\text{keep}=0.2$) & 15.09 & 12.51 & 16.74 & 14.78 \\
DARE ($\text{drop}=0.2$) & 9.80 & 10.00 & 10.00 & 9.93 \\
S-TIES ($\text{keep}=0.2$) & 11.38 & 12.06 & 15.94 & 13.13 \\
S-DARE ($\text{drop}=0.2$) & 9.80 & 10.00 & 10.00 & 9.93 \\
\midrule
\textbf{STDFS (Ours, $\alpha=0.05$)} & 55.90 & 38.55 & \textbf{37.02} & 43.82 \\
\textbf{STDFS (Ours, $\alpha=0.1$)} & 55.00 & \textbf{43.99} & 35.82 & \textbf{44.94} \\
\textbf{STDFS (Ours, $\alpha=0.2$)} & 50.10 & 42.38 & 34.78 & 42.42 \\
\textbf{STDFS (Ours, $\alpha=0.9$)} & 36.59 & 30.85 & 29.37 & 32.27 \\
\midrule
\textbf{SWSM (Ours, $\alpha=0.1, \gamma=0.01$)} & 55.29 & 43.62 & 35.57 & 44.83 \\
\textbf{LASP (Ours, $\beta=0.4$)} & 54.91 & 43.06 & 36.34 & 44.77 \\
\textbf{LASP-Soft (Ours, $\beta=0.4, \gamma=0.02$)} & 55.40 & 42.52 & 36.51 & 44.81 \\
\bottomrule
\end{tabular}
\end{footnotesize}
\end{center}
\vskip -0.1in
\end{table*}

\begin{table*}[t]
\caption{Model merging performance under the \textbf{Shared Head} setting. Values represent test classification accuracy (\%).}
\label{tab:results_shared}
\vskip 0.05in
\begin{center}
\begin{footnotesize}
\begin{tabular}{lcccc}
\toprule
Method & MNIST & Fashion-MNIST & CIFAR-10 & Average \\
\midrule
Weight Averaging (WA) & 30.39 & 27.91 & 26.98 & 28.43 \\
Task Arithmetic (TA, $\lambda=0.3$) & 31.23 & 28.03 & 28.43 & 29.23 \\
\midrule
TIES-Merging ($\text{keep}=0.2$) & 16.54 & 12.41 & 15.55 & 14.83 \\
DARE ($\text{drop}=0.2$) & 9.80 & 10.00 & 10.00 & 9.93 \\
S-TIES ($\text{keep}=0.2$) & 13.29 & 11.95 & 15.41 & 13.55 \\
S-DARE ($\text{drop}=0.2$) & 9.80 & 10.00 & 10.00 & 9.93 \\
\midrule
\textbf{STDFS (Ours, $\alpha=0.1$)} & 53.60 & \textbf{43.44} & 34.60 & 43.88 \\
\midrule
\textbf{SWSM (Ours, $\alpha=0.1, \gamma=0.01$)} & 53.81 & 42.87 & 34.62 & 43.77 \\
\textbf{LASP (Ours, $\beta=0.4$)} & 53.90 & 43.02 & 35.54 & \textbf{44.15} \\
\textbf{LASP-Soft (Ours, $\beta=0.4, \gamma=0.01$)} & 54.02 & 42.67 & 35.67 & 44.12 \\
\textbf{LASP-Soft (Ours, $\beta=0.4, \gamma=0.02$)} & \textbf{54.09} & 42.30 & \textbf{35.76} & 44.05 \\
\bottomrule
\end{tabular}
\end{footnotesize}
\end{center}
\vskip -0.1in
\end{table*}

\begin{figure*}[t]
  \centering
  \begin{subfigure}[b]{0.45\textwidth}
    \centering
    \includegraphics[width=0.95\textwidth]{method_comparison}
    \caption{Performance comparison of model merging methods.}
    \label{fig:method_comparison}
  \end{subfigure}
  \hfill
  \begin{subfigure}[b]{0.45\textwidth}
    \centering
    \includegraphics[width=0.95\textwidth]{alpha_sweep}
    \caption{Impact of the low-frequency ratio $\alpha$ on STDFS.}
    \label{fig:alpha_sweep}
  \end{subfigure}
  \caption{Empirical evaluation of Spectral Task Decoupling in Fourier Space (STDFS) on the vision benchmark. (a) Under the Task-Specific Heads setting, STDFS significantly outperforms standard Weight Averaging, Task Arithmetic, and pruning-based deconfliction methods (TIES, DARE). (b) STDFS performance peaks at $\alpha = 0.1$, proving the need for a frequency-aware balance between shared global structures and task-specific high-frequency details.}
  \label{fig:main_results}
  \vskip -15pt
\end{figure*}

\section{Experimental Setup}
\label{sec:setup}
\subsection{Datasets and Model Architecture}
We evaluate STDFS on a multi-task vision benchmark consisting of three distinct image classification datasets: **MNIST** \cite{lecun1998gradient}, **Fashion-MNIST** \cite{xiao2017fashion}, and **CIFAR-10** \cite{krizhevsky2009learning}. Each dataset contains 10 classes. We use **ResNet-18** \cite{he2016deep} as our core architecture. To ensure compatability, grayscale images from MNIST and Fashion-MNIST are bilinearly resized to $32 \times 32$ pixels and repeated across 3 channels, matching the dimensions of CIFAR-10. All inputs are normalized using standard ImageNet mean and standard deviation.

\subsection{Training and Evaluation Settings}
We initialize our model from ImageNet pre-trained ResNet-18 weights, replacing the classification head with a newly initialized 10-class linear layer. This forms our shared base model $\theta_{\text{base}}$. We then fine-tune $\theta_{\text{base}}$ independently on each task using Adam optimizer ($\text{lr} = 10^{-4}$, batch size 128) to produce three task experts. We evaluate two evaluation settings:
\begin{enumerate}\itemsep=1pt \topsep=1pt \parsep=0pt
    \item \textbf{Task-Specific Heads:} Only the feature extraction backbone is merged. During evaluation on task $k$, we use the merged backbone with the original unmerged classification head of expert $k$.
    \item \textbf{Shared Head:} The entire network, including the classification heads, is merged. A single averaged classification head is used across all tasks.
\end{enumerate}

\section{Results and Analysis}
\label{sec:results}

We present the evaluation results of our proposed STDFS method compared to standard Weight Averaging and Task Arithmetic baselines in \cref{tab:results_task_specific} and \cref{tab:results_shared}.

\subsection{Quantitative Results}
Our proposed STDFS method achieves a remarkable performance improvement over standard Weight Averaging and Task Arithmetic. As shown in \cref{tab:results_task_specific}, Weight Averaging collapses to $30.94\%$ average accuracy under the Task-Specific Heads setting. Task Arithmetic fails to improve beyond Weight Averaging, peaking at $30.91\%$ average accuracy at its optimal $\lambda=0.3$, and completely collapsing to random guessing ($9.93\%$ average accuracy) for larger $\lambda$ due to activation explosion (runaway gradient scales). In contrast, STDFS with a low-frequency ratio of $\alpha = 0.1$ achieves an outstanding **$44.94\%$ average accuracy**, outperforming both standard Weight Averaging and Task Arithmetic by **$+14.0\%$ absolute accuracy**! This dramatic improvement demonstrates that our spectral deconfliction approach successfully resolves the parameter interference bottleneck without experiencing activation scaling instabilities.

Furthermore, under the more challenging \textbf{Shared Head} setting (where the classification heads are also averaged, as shown in \cref{tab:results_shared}), Weight Averaging collapses to $28.43\%$ average accuracy, and Task Arithmetic struggles at $29.23\%$. Our STDFS method with $\alpha = 0.1$ maintains a superb **$43.88\%$ average accuracy**, outperforming standard Weight Averaging by **$+15.45\%$ absolute accuracy** and Task Arithmetic by **$+14.65\%$ absolute accuracy**! This extreme consistency across both head configurations highlights the robust nature of frequency-domain merging.

\subsection{Comparison with Pruning-based Baselines (TIES \& DARE)}
We compare STDFS with two standard parameter-space deconfliction frameworks: TIES-Merging \cite{yadav2023ties} and DARE \cite{yu2024dare}, as well as their frequency-domain equivalents (S-TIES and S-DARE) which we introduce in this work.

As shown in \cref{tab:results_task_specific} and \cref{tab:results_shared}, standard TIES-Merging with a keep ratio of $0.2$ achieves only $14.78\%$ and $14.83\%$ average accuracy under the Task-Specific and Shared Head settings, respectively. DARE with a drop rate of $0.2$ collapses entirely to random guessing ($9.93\%$ average accuracy). Moving these pruning methods to the frequency domain in S-TIES and S-DARE fails to prevent this collapse, with S-TIES and S-DARE yielding $13.13\%$ and $9.93\%$ average accuracy, respectively.

We identify two fundamental reasons why pruning-based model merging is unsuitable for convolutional networks like ResNet-18, contrasting with our spectral routing approach:
\begin{enumerate}\itemsep=1pt \topsep=1pt \parsep=0pt
    \item \textbf{Dense Representation of CNN Backbone:} Unlike Large Language Models (LLMs) or Vision Transformers which possess highly overparameterized, sparse attention heads, the feature representation in convolutional backbones is highly compact and dense. Setting 80\% of weight updates to zero (as in TIES-Merging with $\text{keep}=0.2$) or randomly dropping 20\% of the weights (as in DARE) destroys the spatial feature extraction capabilities of convolutional filters, leading to representational collapse.
    \item \textbf{The Necessity of Global Coarse Structure:} In S-TIES and S-DARE, the pruning is applied uniformly across all frequencies, which zero-outs or drops critical low-frequency coefficients. Since low frequencies capture the fundamental domain-agnostic shapes and templates necessary for routing features, pruning them destroys the model's global coherence.
\end{enumerate}
In contrast, our STDFS framework ensures that 100\% of the low-frequency global structure is preserved and averaged across tasks. Furthermore, instead of zeroing out high-frequency coefficients, STDFS employs magnitude-based routing, which retains the original weight magnitudes and density of all convolutional filters. This frequency-aware treatment explains why STDFS achieves spectacular performance ($44.94\%$), far outperforming all spatial and spectral pruning techniques.

\subsection{Analysis of Low-Frequency Ratio ($\alpha$)}
The parameter $\alpha$ controls the trade-off between feature sharing (averaging low frequencies) and task-specific preservation (routing high frequencies). We observe that a small $\alpha = 0.1$ is optimal. When $\alpha$ is too small (e.g., $0.05$), the shared structural updates are not averaged enough, which slightly degrades MNIST and Fashion-MNIST performance. When $\alpha$ is too large (e.g., $0.9$), too much of the frequency spectrum is averaged linearly, reverting the model toward standard Weight Averaging and re-introducing spatial parameter interference, reducing average accuracy back to $32.27\%$.

\subsection{Empirical Analysis of Layer-wise Spectral Energy Partition}
We validate our core hypothesis that updates partition into coarse, shared low-frequency structures and fine-grained, task-specific details by computing the layer-wise update energy (sum of squared DCT coefficients) across experts with $\alpha=0.1$.

As summarized in \cref{tab:spectral_energy}, we compute the percentage of update energy concentrated in the high-frequency region ($\mathcal{R}_{\text{high}}$) across the network.

\begin{table}[tb]
\caption{Distribution of parameter update energy in the high-frequency region ($\mathcal{R}_{\text{high}}$) across ResNet-18 layer groups with $\alpha = 0.1$. Values represent the percentage (\%) of total squared DCT coefficient magnitude residing in high frequencies.}
\label{tab:spectral_energy}
\vskip 0.05in
\begin{center}
\begin{footnotesize}
\begin{tabular}{lcccc}
\toprule
Layer & MNIST & F-MNIST & CIFAR-10 & Avg. \\
\midrule
\texttt{conv1} & 79.40 & 72.46 & 74.44 & 75.43 \\
\texttt{layer1} & 79.26 & 80.01 & 79.02 & 79.43 \\
\texttt{layer2} & 83.99 & 84.38 & 83.88 & 84.08 \\
\texttt{layer3} & 85.64 & 85.19 & 84.96 & 85.26 \\
\texttt{layer4} & 88.62 & 88.58 & 88.75 & 88.65 \\
\bottomrule
\end{tabular}
\end{footnotesize}
\end{center}
\vskip -0.1in
\end{table}

We observe a clear, monotonic increase in high-frequency energy concentration as depth increases. Early \texttt{conv1} concentrates $75.43\%$ energy in high frequencies, meaning a quarter ($24.57\%$) of its update resides in the $10\%$ low-frequency region. This confirms the domain-agnostic nature of early layers, where coarse templates (e.g., edges, textures) are highly shared. Conversely, deep \texttt{layer4} concentrates $88.65\%$ energy in high frequencies, indicating that deep representation layers (representing semantic abstractions and classification boundaries) require highly localized, fine-grained adjustments.

This empirical finding rigorously justifies STDFS: (1) low-frequency linear averaging is optimal in early layers due to high template correlation; (2) task-specific details are concentrated in high frequencies in deep layers, where linear averaging blurs coefficients and causes activation variance collapse, whereas spectral max-pooling preserves update scale.

\subsection{Evaluation of Layer-wise Adaptive Spectral Partitioning (LASP)}
We evaluate Layer-wise Adaptive Spectral Partitioning (LASP) under both configurations. Under Task-Specific Heads (\cref{tab:results_task_specific}), LASP ($\beta=0.4$) achieves \textbf{44.77\%} average accuracy. Crucially, on CIFAR-10, LASP scores \textbf{36.34\%}, a substantial \textbf{+0.52\%} improvement over uniform STDFS (35.82\%). 

Under Shared Heads (\cref{tab:results_shared}), LASP ($\beta=0.4$) achieves a superb \textbf{44.15\%} average accuracy, outperforming uniform STDFS (43.88\%) by \textbf{+0.27\%} and Weight Averaging by \textbf{+15.72\%}. On CIFAR-10, LASP scores \textbf{35.54\%}, beating STDFS by a massive \textbf{+0.94\%}.

These results validate the core philosophy of LASP: matching the low-frequency partitioning threshold $\alpha^l$ to the layer-wise empirical energy profile allows earlier layers to perform wider feature sharing, while restricting deeper layers to highly selective spectral max-pooling. This layer-aware routing preserves localized high-frequency semantic features in deeper representation layers, yielding superior classification.

\subsection{Evaluation of Soft Boundary Transitions (SWSM \& LASP-Soft)}
Traditional STDFS uses a hard step-function cutoff, which can introduce frequency-domain boundary discontinuities (similar to Gibbs ringing artifacts). We propose Spectral Window-based Soft Merging (SWSM) and Layer-wise Adaptive Soft Spectral Partitioning (LASP-Soft), which use a sigmoid gating function with temperature $\gamma$ to smoothly blend averaged low-frequency and max-pooled high-frequency updates.

Sweeping $\gamma \in \{0.01, \dots, 0.2\}$ under both head settings (\cref{tab:results_task_specific}, \cref{tab:results_shared}) yields two core insights. First, the partition between domain-agnostic and task-specific parameters is remarkably sharp: under SWSM with a tight window ($\gamma=0.01$), average accuracy is \textbf{44.83\%} (highly consistent with STDFS's 44.94\%), but as the transition window widens ($\gamma \ge 0.1$), performance degrades monotonically. Second, smoothly blending the boundary resolves Gibbs artifacts on complex semantics: on CIFAR-10, LASP-Soft ($\gamma=0.02$) scores a spectacular \textbf{36.51\%} (Task-Specific) and \textbf{35.76\%} (Shared), outperforming hard LASP (+0.17\% / +0.22\%) and uniform STDFS (+0.69\% / +1.16\%). For semantically dense tasks requiring filter continuity, a soft spectral transition window effectively suppresses boundary artifacts, enabling experts to preserve fine-grained representations.

\section{Conclusion}
\label{sec:conclusion}
We introduce **Spectral Task Decoupling in Fourier Space (STDFS)**, a training-free and data-free model merging framework operating in the frequency domain. STDFS applies a 2D-DCT to weight updates, averaging low-frequency structural features while routing high-frequency task-specific details via spectral max-pooling. On a diverse vision benchmark, STDFS outperforms standard Weight Averaging by over $+14\%$ absolute accuracy. Our extensions, LASP and LASP-Soft, further improve performance by dynamically adapting frequency boundaries and smoothing spectral transitions. This work establishes frequency-domain deconfliction as a powerful paradigm for parameter fusion, opening exciting avenues for model merging.

\bibliography{example_paper}
\bibliographystyle{icml2026}

%%%%%%%%%%%%%%%%%%%%%%%%%%%%%%%%%%%%%%%%%%%%%%%%%%%%%%%%%%%%%%%%%%%%%%%%%%%%%%%
%%%%%%%%%%%%%%%%%%%%%%%%%%%%%%%%%%%%%%%%%%%%%%%%%%%%%%%%%%%%%%%%%%%%%%%%%%%%%%%
% APPENDIX
%%%%%%%%%%%%%%%%%%%%%%%%%%%%%%%%%%%%%%%%%%%%%%%%%%%%%%%%%%%%%%%%%%%%%%%%%%%%%%%
%%%%%%%%%%%%%%%%%%%%%%%%%%%%%%%%%%%%%%%%%%%%%%%%%%%%%%%%%%%%%%%%%%%%%%%%%%%%%%%
\clearpage
\appendix
\onecolumn
\section{Spectral Task Alignment Analysis}
\label{app:spectral_similarity}

To further validate the core hypothesis of Spectral Task Decoupling in Fourier Space (STDFS), we investigate the geometric alignment (cosine similarity) of task-specific parameter updates (task vectors) in both the spatial and spectral domains. Specifically, for any pair of tasks $i$ and $j$, and a given weight matrix update $U_i, U_j$, we compute:
\begin{enumerate}
    \item \textbf{Spatial Cosine Similarity:} 
    \[
    \cos(U_i, U_j) = \frac{\langle U_i, U_j \rangle}{\|U_i\|_2 \|U_j\|_2}
    \]
    \item \textbf{Low-Frequency Spectral Similarity:} 
    \[
    \cos_{\text{low}}(U_i, U_j) = \frac{\langle \mathcal{D}_i \odot \mathcal{M}_{\text{low}}, \mathcal{D}_j \odot \mathcal{M}_{\text{low}} \rangle}{\|\mathcal{D}_i \odot \mathcal{M}_{\text{low}}\|_2 \|\mathcal{D}_j \odot \mathcal{M}_{\text{low}}\|_2}
    \]
    \item \textbf{High-Frequency Spectral Similarity:} 
    \[
    \cos_{\text{high}}(U_i, U_j) = \frac{\langle \mathcal{D}_i \odot \mathcal{M}_{\text{high}}, \mathcal{D}_j \odot \mathcal{M}_{\text{high}} \rangle}{\|\mathcal{D}_i \odot \mathcal{M}_{\text{high}}\|_2 \|\mathcal{D}_j \odot \mathcal{M}_{\text{high}}\|_2}
    \]
\end{enumerate}
where $\mathcal{D}_k = \text{2D-DCT}(U_k)$, and $\mathcal{M}_{\text{low}}, \mathcal{M}_{\text{high}}$ represent the binary partitioning masks defined by the low-frequency ratio $\alpha = 0.1$. 

We compute these metrics across all convolutional layer groups and average them over all pairwise task combinations (MNIST vs. Fashion-MNIST, MNIST vs. CIFAR-10, and Fashion-MNIST vs. CIFAR-10). The empirical results are visualized in \cref{fig:spectral_similarity}.

\begin{figure}[ht]
\centering
\includegraphics[width=0.6\textwidth]{spectral_similarity_comparison.pdf}
\caption{Average cosine similarity between task updates across ResNet-18 layer groups in the spatial domain vs. the low-frequency and high-frequency spectral bands ($\alpha = 0.1$). Although task updates are highly orthogonal overall, low-frequency similarity is consistently higher, confirming that shared structural updates reside in the low-frequency spectrum, while high-frequency features are highly task-specific and orthogonal.}
\label{fig:spectral_similarity}
\end{figure}

\subsection{Key Findings and Insights}
\begin{itemize}
    \item \textbf{Near-Perfect Orthogonality in the Spatial Domain:} The overall average spatial cosine similarity across all layer groups is extremely low ($0.0040$), indicating that the fine-tuning updates for different tasks starting from a shared pre-trained initialization are highly orthogonal. This is expected given the extreme divergence of the input domains and modalities (e.g., grayscale digits vs. color natural images).
    \item \textbf{Concentration of Alignment in Low Frequencies:} Despite this near-perfect spatial orthogonality, the average similarity in the low-frequency band is $0.0148$, which is nearly \textbf{4x larger} than the spatial similarity ($0.0040$), and nearly \textbf{15x larger} than the high-frequency similarity ($0.0010$). This represents a major empirical confirmation of our core hypothesis: any commonalities or shared representation structures that exist between disparate tasks are heavily concentrated in the low-frequency coefficients of the parameter updates.
    \item \textbf{High-Frequency Divergence:} In contrast, the high-frequency similarity decays almost to absolute orthogonality ($0.0010$ overall, and even negative in early layers like \texttt{conv1}). This demonstrates that the fine-grained details where task conflicts primarily occur are completely uncorrelated.
\end{itemize}

These findings provide a powerful, mathematically rigorous justification for the design of STDFS:
\begin{enumerate}
    \item Since low-frequency components exhibit higher correlation and capture shared structural templates, linear averaging (Weight Averaging) is both safe and beneficial in this band, helping consolidate domain-agnostic capabilities.
    \item Since high-frequency components are highly orthogonal and task-specific, linear averaging blurs these distinct coefficients, collapsing their magnitude and causing activation variance decay. By applying relative spectral max-pooling instead of averaging, STDFS selects the dominant task-specific filters without destructive interference, preserving parameter energy and model capacity.
\end{enumerate}

\end{document}
